\chapter{Bakgrund och teori}
Detta kapitel presenterar den teoretiska bakgrund som krävs för att förstå social engineering som ett cybersäkerhetshot. Inledningsvis diskuteras cybersäkerhetens roll i moderna digitala system och grunden för att förstå cybersäkerhet i dagens digitala miljöer. Kapitlet kommer också att behandla social engineering som ett fenomen, med fokus på dess unika angreppssätt som utnyttjar mänskliga faktorer snarare än enbart tekniska sårbarheter. Avslutningsvis diskuteras de psykologiska mekanismer som ligger till grund för social engineering-attacker och en genomgång av vanliga angreppsmetoder och attacktyper.

\section{Cybersäkerhet i moderna digitala system}
Under de senaste årtiondena har digitaliseringen lett till att samhällsfunktioner, organisationer och individer blivit alltmer beroende av digitala system. Utvecklingen har möjliggjort effektivare verksamhetsprocesser och skapat ett överflöd av nya digitala tjänster. Däremot har utvecklingen också samtidigt ökat systemens exponering för cyberhot och expanderat attackytan för hotaktörer \cite{gururaj2024social_engineering}.

För moderna digitala system kan cybersäkerhet förstås som skyddet av digital information och system genom tekniska, organisatoriska och mänskliga åtgärder. Enligt \cite{washo2021interdisciplinary, enisa2018cybersecurity_culture} ingår det i cybersäkerhet inte enbart teknik, utan också socio-tekniska fenomen där mänskliga faktorer spelar en central roll. Detta innebär att ett ensidigt fokus på tekniska sårbarheter inte är tillräckligt för att skydda moderna digitala system.

Därtill kännetecknas dagens digitala system av en hög grad av sammankoppling och komplexitet. Digitala system består av flera lager, tjänster och aktörer som samverkar i nätverksbaserade strukturer \cite{springer_social_engineering_ontology}. Enligt \cite{springer_social_engineering_ontology} innebär strukturen att sårbarheter i en del av systemet kan leda till konsekvenser i andra delar. Komplexiteten leder till försvårad riskbedömning, säkerhetsstyrning och incidenthantering. 

Därpå har den ökade användningen av digital kommunikation, såsom e-post, nätbaserade plattformar och e-tjänster, skapat nya möjligheter för angripare att nå målgrupper med begränsade resurser \cite{basnet2023anti_phishing}. Även om tekniska försvarsmekanismer utvecklas kontinuerligt, anpassar angripare sina metoder i takt med förbättrad cybersäkerhet. 

För att sammanfatta kännetecknas moderna digitala system av komplexitet, sammankoppling och ett beroende av mänsklig interaktion. Enligt \cite{washo2021interdisciplinary, enisa2018cybersecurity_culture} kräver cybersäkerheten i dessa miljöer ett integrerat perspektiv där tekniska lösningar kombineras med förståelse för organisatoriska och mänskliga faktorer.

\section{Social Engineering som cybersäkerhetshot}
Begreppet social engineering har historiskt använts inom samhällsvetenskap för att beskriva planerade försök att forma sociala strukturer och mänskligt beteende. I modern cybersäkerhetskontext har begreppet däremot fått en mer begransad betydelse och syftar främst på manipulation av individer i digitala miljöer \cite{podgorecki1998social}.

Social engineering utgör en attackform där tekniska verktyg och psykologisk manipulation kombineras för att påverka individers beteende i digitala miljöer \cite{springer_social_engineering_ontology, gururaj2024social_engineering}. Till skillnad från traditionella cyberhot, som fokuserar på rent tekniska intrång, strävar social engineering efter att få användare att frivilligt kringgå etablerade säkerhetsåtgärder. Hotaktören manipulerar därmed den mänskliga länken i säkerhetskedjan snarare än den tekniska infrastrukturen i sig.

Social engineering kan därmed kännetecknas av att den utnyttjar kommunikationskanaler och tekniska plattformar för att skapa trovärdiga scenarier \cite{se_threat_defense_survey, algarni2019se_survey}. Kommunikationskanalerna och plattformarna fungerar här som möjliggörare: e-post, sociala medier och andra digitala system används för att distribuera budskap utformat som legitima meddelanden. Dessa meddelanden är dock färdigt konstruerade fällor, där tekniskt byggda url-länkar kan stjäla information eller leda till andra cyberhot. Attacken bygger samtidigt på psykologiska mekanismer såsom auktoritet, brådska eller tillit, vilket ökar sannolikheten för att offret agerar utan att ifrågasätta situationen.

Kombinationen av teknisk skicklighet och psykologisk påverkan gör social engineering till ett särskilt effektivt cyberhot. Även i tekniskt välskyddade miljöer kan en användare, genom manipulation, bli den punkt där säkerheten bryts \cite{washo2021interdisciplinary}. Social engineering kan därmed fungera som en bro mellan mänskligt beteende och tekniskt intrång, vilket är orsaken till att den ofta utgör den initiala fasen i mer omfattande cyberattacker \cite{algarni2019se_survey}. 

Vidare visar det sig att social engineering inte är en enskild metod för cyberattacker, utan en samling attackstrategier som kan anpassas av hotaktören enligt kontext och mål \cite{aldawood2019taxonomy, gururaj2024social_engineering}. Detta pekar på att hotet är genomtänkt och välstrukturerat istället för godtyckligt. Metodens flexibilitet, i kombination med beroende av både tekniska och psykologiska faktorer, gör den svår att motverka genom allmänna skyddsåtgärder \cite{se_threat_defense_survey}.

\section{Psykologiska mekanismer bakom social engineering}
Som tidigare nämnt bygger social engineering i grunden på psykologisk påverkan snarare än tekniskt utnyttjande. Angreppsmetodens effektivitet förklaras av dess utnyttjande av återkommande mönster i mänskligt beslutsfattande och social interaktion \cite{psychology_social_engineering}. I stället för att kringgå tekniska säkerhetssystem påverkas individen direkt att själv utföra den handling som möjliggör intrånget. 

Den underliggande mekanismen som ligger bakom social engineering är dess användning av kognitiva biaser, det vill säga systematiska tankefel som påverkar hur människor bedömer information och risk \cite{psychology_social_engineering}. Som exempel kan auktoritetsbias leda till att en individ, utan närmare granskning, följer instruktioner från någon som uppfattas ha en chefsposition. På liknande sätt kan tidspress eller känslomässig påverkan minska benägenheten att ifrågasätta en viss situation. 

Utöver individuella beslutsprocesser utnyttjar social engineering även sociala normer och organisatoriska strukturer \cite{washo2021interdisciplinary}. Samarbetsvilja, tillit och respekt för hierarkier är nödvändiga för att organisationer ska fungera effektivt, men dessa faktorer kan samtidigt missbrukas av hotaktörer. Användningen av social engineering mot organisationer blir därmed möjlig eftersom de normer som tillåter fungerande arbetsmiljöer också skapar förutsättningar för manipulation.


Enligt översiktsstudier visar det sig att användare ofta underskattar sannolikheten att själva bli utsatta för manipulation \cite{algarni2019se_survey}. Denna övertro på den egna förmågan att identifiera hot bidrar oftast till att enkla säkerhetsrutiner kringgås, även när medvetenhet om dessa risker existerar. Den psykologiska delen av social engineering riktar sig däremot inte enbart mot okunskap, utan även mot övertro och rutinmässigt beteende. 

Enligt Europeiska unionens cybersäkerhetsbyrå (ENISA) påverkas mänskligt beteende i säkerhetskontext av både individuella och organisatoriska faktorer \cite{enisa2018cybersecurity_culture}. Arbetsbelastning, säkerhetskultur och tydlighet i rutiner spelar en stor roll för hur individer reagerar på potentiella hot. När dessa faktorer samverkar, leder det till att social engineering oftast lyckas och gör det svårt för individer att förstå att de har blivit utsatta för manipulation. 

\section{Angreppsmetoder och attacktyper}
Social engineering omfattar en stor mängd olika angreppsmetoder som kan kategoriseras utifrån angreppsätt, kommunikationskanal eller tekniskt genomförande. För att tydliggöra hur manipulationen genomförs i praktiken har dessa angreppsmetoder strukturerats i taxonomiska modeller \cite{aldawood2019taxonomy, algarni2024comprehensive_taxonomy, gururaj2024social_engineering, krombholz2015overview}. 

En av de mest allmänna attacktyperna är nätfiske, där hotaktören via e-post, meddelandetjänster eller webbplatser försöker lura mottagaren att lämna ut känslig information eller klicka på skadliga länkar \cite{se_threat_defense_survey, algarni2019se_survey}. En mer riktad variant av nätfiske kallas för spjutfiske, där attacken anpassas efter en specifik individ eller organisation genom insamlad bakgrundsinformation \cite{heartfield_loukas_se}. Denna form av nätfiske anses ofta vara mer framgångsrik eftersom attacken upplevs som legitim och trovärdig.

Utöver e-postbaserade angrepp förekommer också pretexting som ett attackhot, där hotaktören skapar ett trovärdigt scenario eller berättelse för att manipulera offret att lämna ut känslig information \cite{mouton_se_systematic}. Målgruppen för pretextingbaserade attacker är oftast organisatoriska roller, där hotaktören exempelvis låtsas vara IT-personal eller en överordnad chef. Andra vanliga attacktyper inkluderar också baiting, där offret manipuleras med lockande erbjudanden eller belöningar, och quid pro quo, där hotaktören erbjuder en tjänst i utbyte mot information \cite{se_threat_defense_survey, algarni2019se_survey}. 

I sociala medier förekommer även attacker där förfalskade profiler används för att bygga upp förtroende över tid innan manipulationen genomförs \cite{fire2021se_social_networks}. Enligt den taxonomiska studien \cite{algarni2024comprehensive_taxonomy} kombinerar moderna social engineering-attacker ofta flera tekniker och kan integrera tekniska komponenter såsom domänförfalskning eller automatiserade meddelandesystem. Detta innebär att gränsen mellan tekniska och mänskliga angrepp blir alltmer otydlig. 


\subsection{Nätfiske och spjutfiske}
Som tidigare nämnts bygger nätfiskeattacker på att hotaktörer utger sig för att representera en legitim organisation eller tjänst i syfte att lura mottagaren att lämna ut känslig information \cite{se_threat_defense_survey, algarni2019se_survey, krombholz2015overview}. Information som hotaktörer oftast är ute efter kan inkludera inloggningsuppgifter, kreditkortsdata eller andra konfidentiella data. 

Attackerna distribueras i de flesta fall via e-post, men förekommer även genom SMS (smishing), sociala medier och förfalskade webbplatser \cite{algarni2019se_survey}. Den tekniska komponenten i nätfiskeattacker består ofta av e-postspoofing, manipulerade URL:er, domänförfalskning eller visuellt övertygande kopior av autentiska inloggningssidor \cite{algarni2019se_survey, basnet2023anti_phishing}. Attackens effektivitet uppstår genom att individer blir lurade av teknik och psykologisk påverkan för att uppfatta budskapet som äkta.

För att nätfiske ska lyckas krävs det en betydande mängd psykologisk manipulation. Enligt \cite{psychology_social_engineering} utnyttjar hotaktörer ofta auktoritet, brådska och rädsla för att minska individers kritiska tänkande. Ett exempel är meddelanden som antyder att ett bankkonto äventyras, vilket skapar ett stressmoment som kan leda till impulsiva beslut.

Till skillnad från traditionella nätfiskeattacker, som massutskickar attacker oriktat, riktar sig spjutfiske mot en specifik organisation eller individ \cite{heartfield_loukas_se}. Spjutfiske kan även anta formen av så kallad whaling, där högt uppsatta personer inom organisationer utgör målgruppen. Attacken bygger på insamlad bakgrundsinformation om offret, exempelvis genom öppna källor eller sociala nätverk. Den insamlade informationen möjliggör personligt anpassade meddelanden som ökar trovärdigheten och sannolikheten för en lyckad attack. 

Moderna nätfiske- och spjutfiskeattacker integrerar allt fler tekniska och sociala komponenter \cite{algarni2024comprehensive_taxonomy}. Exempelvis kan hotaktörer kombinera en riktad informationsinsamling med artificiell intelligens (AI) för att generera språkligt sofistikerade meddelanden \cite{basnet2023anti_phishing}. Forskning visar att användningen av AI i nätfiske- eller spjutfiskesammanhang har ökat attackernas trovärdighet och därmed försvårat åtskiljandet mellan legitim och bedräglig kommunikation \cite{zhang2024digital_deception}.

\subsection{Pretexting och identitetsbedrägeri}
Till skillnad från nätfiske, där attacken ofta baseras på massdistribution, utgör pretexting en strukturerad och direkt form av social engineering. Hotaktören skapar ett trovärdigt scenario eller antar en falsk identitet i syfte att manipulera offret \cite{se_threat_defense_survey, zhang2024digital_deception, krombholz2015overview}. 

Grunden i pretexting är skapandet av ett sammanhang där hotaktören framstår som legitim. Beroende på målgruppen kan hotaktören exempelvis utge sig för att vara en arbetssökande, IT-personal eller en annan intern funktion inom en organisation \cite{algarni2019se_survey}. Genom att utnyttja etablerade organisatoriska roller och hierarkier kan hotaktören minska sannolikheten för att manipulationen ifrågasätts.

Identitetsbedrägeri är nära kopplat till pretexting och innebär att hotaktörer aktivt imiterar en verklig person eller institution. Imitation kan exempelvis ske genom falska e-postadresser, manipulerade telefonnummer (caller ID spoofing) eller förfalskade profiler på sociala medier \cite{aldawood2019taxonomy}. Även här kombineras tekniska verktyg med social manipulation för att skapa en övertygande illusion av autenticitet. 

Enligt \cite{washo2021interdisciplinary} utnyttjar  pretextingbaserade attacker organisatoriska normer såsom samarbetsvilja, respekt för auktoritet och hjälpsamhet. Normerna är nödvändiga för att organisationer ska fungera effektivt, men de kan samtidigt utnyttjas av potentiella hotaktörer. Orsaken till att pretexting oftast ingår i bredare social engineering-attacker är att attacken inte enbart fokuserar på individers okunskap, utan också på strukturella och kulturella faktorer inom organisationer.

\subsection{Insider-relaterade attacker}
Insider-relaterade attacker skiljer sig från tidigare diskuterade attacker genom att de involverar individer som redan har tillgång till en organisations system från insidan \cite{algarni2019se_survey}. Till skillnad från externa attacker bygger dessa hot, som sagt, på att en intern aktör, medvetet eller omedvetet, bidrar till attacken. Insiderhot kan delas in i två huvudkategorier: avsiktliga och oavsiktliga. Avsiktliga insiderattacker innebär att en anställd medvetet missbrukar sina åtkomsträttigheter, exempelvis för ekonomisk vinst eller sabotage \cite{algarni2024comprehensive_taxonomy, gururaj2024social_engineering}. Oavsiktliga insiderincidenter uppstår däremot när en anställd manipuleras genom social engineering och därmed omedvetet möjliggör ett intrång. 

Som redan nämnts utgör social engineering ofta en utgångspunkt för insiderrelaterade incidenter. Genom nätfiske, spjutfiske eller pretexting kan hotaktören påverka interna användare inom en organisation att kringgå säkerhetsrutiner eller ge obehörig åtkomst till system \cite{algarni2019se_survey}. Forskning visar att bristande medvetenhet eller otydlig ansvarsfördelning kan skapa situationer där manipulation är mer framgångsrik. Därför spelar organisatoriska faktorer såsom säkerhetskultur, utbildningsnivå och tydliga rutiner en avgörande roll i hur mottaglig en organisation är för insider-relaterade attacker \cite{albladi2018defense_mechanisms, enisa2018cybersecurity_culture}.